\newcommand{\gptpart}{30}
\newcommand{\gpttitle}{Conclusion}
\newcommand{\gptsubtitle}{The Restoration of Reality}
\newcommand{\gptdate}{November 2026}
% Shared preamble for the second edition of Graviton Pressure Theory.
% A part's own file sets \gptpart, \gpttitle, \gptsubtitle and \gptdate, then \input{../common/preamble}.
\documentclass[11pt,a4paper]{article}
\usepackage[T1]{fontenc}
\usepackage[utf8]{inputenc}
\usepackage{mathpazo}
\usepackage[scaled=0.92]{helvet}
\usepackage{amsmath,amssymb}
\usepackage[margin=2.5cm]{geometry}
\usepackage{microtype}
\usepackage{graphicx}
\usepackage{booktabs,array,tabularx}
\usepackage[font=small,labelfont=bf]{caption}
\usepackage{enumitem}
\usepackage[authoryear,round]{natbib}
\usepackage{titlesec}
\usepackage[titles]{tocloft}
\renewcommand{\cftsecfont}{\bfseries}\renewcommand{\cftsecpagefont}{\bfseries}
\renewcommand{\cftsecleader}{\hfill}\setlength{\cftsecnumwidth}{3.2em}
\usepackage{gpt-status}
\usepackage[hidelinks]{hyperref}
\setlength{\parskip}{4pt}
\setlength{\emergencystretch}{3em}
\setlength{\parindent}{0pt}
\setlist{itemsep=2pt,topsep=3pt}
\titleformat{\section}{\Large\bfseries}{\thesection}{1em}{}
\titleformat{\subsection}{\large\bfseries}{\thesubsection}{1em}{}
\renewcommand{\thesection}{\gptpart.\arabic{section}}
\renewcommand{\thetable}{\gptpart.\arabic{table}}
\renewcommand{\thefigure}{\gptpart.\arabic{figure}}
\renewcommand{\theequation}{\gptpart.\arabic{equation}}
\setcounter{tocdepth}{1}
\makeatletter\renewcommand{\@pnumwidth}{2.2em}\renewcommand{\@tocrmarg}{3em}
\makeatother
\bibliographystyle{plainnat}
\newcommand{\gptitle}{%
  \begin{titlepage}\centering
  \vspace*{3cm}
  {\Huge\bfseries Graviton Pressure Theory\par}\vspace{0.8cm}
  {\Large The Unified Framework\par}\vspace{0.2cm}
  {\large Individual Submission\par}\vspace{0.6cm}
  {\normalsize This document is part of a multi-part scientific framework\par}\vspace{2cm}
  {\normalsize Part \gptpart\ of 30\par}\vspace{0.8cm}
  {\huge\bfseries \gpttitle\par}\vspace{0.6cm}
  {\Large \gptsubtitle\par}\vspace{1.5cm}
  {\small\begin{minipage}{0.75\textwidth}\centering This submission is part of the broader Graviton Pressure Theory (GPT) project, a redefinition of gravitational interaction rooted in causal field dynamics and physical force transmission. Each document is designed to stand on its own; its full context emerges as part of the larger framework. For the complete series, see https://peopleoftheworld.net/graviton-pressure-theory/.\end{minipage}\par}
  \vfill
  {\large Author: Shareef Ali Rashada\par}\vspace{0.2cm}
  {\normalsize shareef@peopleoftheworld.net\par}\vspace{0.5cm}
  {\normalsize Second edition, \gptdate \quad (first edition 12 June 2025)\par}
  \end{titlepage}
  \setcounter{page}{1}\tableofcontents\clearpage}

\graphicspath{{figs/}}
\begin{document}
\gptitle

\section*{About This Paper}
\addcontentsline{toc}{section}{About This Paper}

Twenty-nine parts have built Graviton Pressure Theory from its question to its tests. This last part says, in one place, what the framework claims, what it has shown, what it has withdrawn, what it still owes, and how it can be proved wrong.

This is the second edition; what changed from the first is listed at the end. Claims are marked by their status, as the key shows.

\statuskey

\section{The Question Newton Left Open}

In 1693 Isaac Newton wrote that a body acting on another at a distance, through nothing, was to him ``so great an absurdity'' that no competent thinker could fall into it.\fn{newtonBentley1693} He had given the law of gravity, and he left its cause open.

Physics since has answered the question \emph{how much} with ever greater precision, and set aside the question \emph{what}. General Relativity gives every number gravity has been asked for: the bending of light, the slowing of clocks, the turning of Mercury's orbit, the waves from merging black holes. It does so by describing gravity as the geometry of spacetime, and by saying that a falling body has nothing acting on it at all (Part 8).\fn{rashadaPart8}

\argm{Graviton Pressure Theory begins where Newton stopped.} It asks what touches what. Its answer is a medium: a substance that fills space, flows into every mass, and presses on everything it meets. Einstein himself said in 1920 that under General Relativity ``space without ether is unthinkable''.\fn{einstein1920} The question of the medium was set aside, not answered. This framework answers it.

\section{One Medium, One Inflow}\label{sec:one}

\hypm{Space is filled with a medium of gravitons, and every mass takes it in.} The medium flows inward toward a mass at the speed $v=\sqrt{2GM/r}$, 11.2 kilometres a second at the surface of the Earth, and its pressure falls as it flows, $P/P_0 = 1 - v^2/c^2$ (Parts 5, 14 and 23).\fn{rashadaPart5}\textsuperscript{,}\fn{rashadaPart14}\textsuperscript{,}\fn{rashadaPart23} \argm{From that one inflow, and nothing fitted, the numbers follow} (Figure~\ref{fig:one}):

\begin{figure}[t]\centering
\includegraphics[width=0.86\linewidth]{fig30-1-one-medium.pdf}
\caption{One medium, one inflow. In blue, measured results the inflow reproduces with the standard numbers; in green, a measurement that supports it; in orange, predictions that are untested; in grey, what is still owed.}\label{fig:one}
\end{figure}

\begin{itemize}
\item \textbf{Falling.} The inflow delivers momentum to every particle of a body in proportion to its mass, so every body falls alike, as measured to one part in $10^{15}$.\fn{touboul2022} Its energy of motion is paid by the medium's pressure turning into flow: the energy of gravity has a place (Parts 7 and 11).\fn{rashadaPart7}\textsuperscript{,}\fn{rashadaPart11}
\item \textbf{Weight.} Held, the same delivery passes through the body into its support, and every scale reading measures it. Holding is throughput, not consumption: the medium's momentum passes through and is not used up (Part 6).\fn{rashadaPart6}
\item \textbf{Clocks.} A clock is a structure of the medium's waves, slowed by the medium streaming past it. Lower clocks run slow, and moving clocks run slow, by one effect: GPS satellite clocks gain 38.5 microseconds a day (Parts 9 and 18).\fn{rashadaPart9}\textsuperscript{,}\fn{rashadaPart18}
\item \textbf{Light.} Light is a wave of the medium, refracted where its pressure is lowered: bent 1.75 arcseconds at the Sun's edge, and delayed passing it (Parts 23 and 28).\fn{rashadaPart28}
\item \textbf{Orbits.} The lowered pressure near the Sun turns Mercury's orbit by 43 arcseconds a century, beyond the planets' pulls (Part 26).\fn{rashadaPart26}
\item \textbf{Spin.} A spinning mass winds the medium around it; the swirl turned Gravity Probe B's gyroscopes by 37 milliarcseconds a year (Part 16).\fn{rashadaPart16}\textsuperscript{,}\fn{everitt2011}
\item \textbf{Waves.} Gravitational waves are transverse waves of the medium and travel at its wave speed, the speed of light, as measured to a few parts in $10^{15}$ (Part 9).\fn{abbott2017speed}
\item \textbf{Horizons.} Where the inflow reaches the speed of light, at $r=2GM/c^2$, light can no longer climb out: a black hole's horizon (Part 25).\fn{rashadaPart25}
\item \textbf{Galaxies.} Far from a mass, where the inflow's acceleration falls to $a_0\approx1.2\times10^{-10}$ m/s$^2$, the medium is brought back into step and gravity changes its form. Galaxies follow one relation between their visible matter and their rotation, with the change at $a_0$ (Part 9).\fn{mcgaugh2016} \hypm{GPT reads it as the relaxation boundary,} where the standard account adds unseen matter.
\end{itemize}

\hypm{And the same medium builds what it acts on.} Matter is energy slowed down: circulation at the speed of light, held in a knot of the medium and fed by its inflow (Parts 17 and 19).\fn{rashadaPart17}\textsuperscript{,}\fn{rashadaPart19} Time is the rate at which a structure renews itself (Part 18). Order comes in levels, built as the universe cooled and removed by heat in reverse (Part 20).\fn{rashadaPart20} Load, heat and charge flow through matter along paths that its order sets (Part 21).\fn{rashadaPart21} A magnet is the medium's swirl, organised and held (Part 22).\fn{rashadaPart22} Light, sound and electricity are waves and flows in one medium (Part 28). Every force names its carrier (Part 27).\fn{rashadaPart27}

\section{What Thirty Parts Have Built}

\begin{table}[t]\centering\footnotesize
\caption{The framework, part by part.}
\renewcommand{\arraystretch}{1.1}\begin{tabularx}{\linewidth}{@{}>{\raggedright\arraybackslash}p{0.13\linewidth} >{\raggedright\arraybackslash}X@{}}\toprule
\textbf{Parts} & \textbf{What they do}\\\midrule
1--4 & State the question: what gravity does, how it has been explained, how description came to stand in for mechanism, and what still does not fit\\
5--8 & Set out the mechanism: the inflow, holding as throughput, every force with its carrier, falling and holding as different acts; and the break with General Relativity, stated fairly\\
9--13 & Weigh the evidence, grade it, and test the mechanism against the four forces, the limits of knowledge, and the tides and rhythms of life\\
14--19 & Build the foundations: the framework in one piece, the graviton, the field's inflow and swirl, mass as circulation, time as renewal, and the pathway to matter\\
20--22 & Order within matter: coherence, the corridors along which flows pass, and magnetism\\
23--28 & The numbers: the equations and what they reproduce, units and constants, the inflow to the speed of light, orbits, every force and its carrier, light and the waves of the medium\\
29 & Every result read as a medium, what a medium must survive, and the tests that could decide\\\bottomrule
\end{tabularx}\label{tab:parts}\end{table}


Table~\ref{tab:parts} gathers the thirty parts. \argm{Read together, they give one account:} one medium, one inflow, one rule for what pushes and what holds, and the same few quantities, the inflow's speed and the medium's pressure, from the fall of an apple to the edge of a black hole.

\section{The Break, Stated Exactly}

\argm{General Relativity and GPT can agree on every number.} In 1921 and 1922 Painlevé and Gullstrand wrote General Relativity's field outside a mass as flat space through which something flows inward at the escape speed, the same $\sqrt{2GM/r}$.\fn{painleve1921}\textsuperscript{,}\fn{gullstrand1922} One set of equations carries both readings, the curved geometry and the flowing medium, and no measurement of a number can choose between them (Parts 8 and 29).\fn{rashadaPart29}

\argm{They cannot agree on what is there.} General Relativity says gravity is not a force, that nothing acts on a falling body, and that gravity's energy cannot be placed at any point.\fn{misner1973} GPT says a medium is there; that it delivers momentum to every body it carries; and that the energy of gravity is in its pressure. These are not two versions of one claim. On what exists, only one can be right.

\argm{The choice rests on mechanism, and here the medium names what the geometry leaves unnamed:} what acts on the falling body, what slows the clock, what slows the light, what turns the gyroscope, what carries a wave and its energy. A geometry gives the path. A medium gives the cause. GPT keeps General Relativity's numbers and supplies the ground they never had.

\section{What Was Withdrawn, and Why It Matters}

\argm{A framework is measured by what it will give up.} In this edition, every claim was set against the measurements, and what failed was withdrawn or restated. Among them:
\begin{itemize}
\item that holding a weight consumes energy (restated: holding still requires energy, a continuous flow through the load that passes through and is not used up, Part 6);
\item that a body's weight depends on its surface, or that corridors change its weight in a fall (restated: the delivery tracks a body's energy, and every body falls alike to one part in $10^{15}$; whether a held body's weight depends on its structure is open, Parts 15, 20 and 21);
\item that heating adds mass beyond $\Delta E/c^2$, or that matter measurably shields gravity (bounded by measurement, Parts 15 and 17);
\item that demonstrations in the classroom validate GPT (restated: most show the same outcome under both accounts, and three point to real tests, Part 29);
\item claims of applications, shielding as engineering, and speculation far from the core.
\end{itemize}
What remains is smaller, and harder. Every claim in it is marked for what it is: argued, a hypothesis, or open.

\section{What Is Still Owed}

\openm{The framework owes these, and says so:}
\begin{itemize}
\item galaxy clusters, which need about twice their visible mass, and the third peak of the background light (Part 9);
\item the law beyond the relaxation boundary, and the value of $a_0$ from the rate at which the medium is brought back into step;
\item how a pair of stars' boundary behaves inside the galaxy's own flow (Part 29);
\item the aberration of gravity from GPT's own inflow, which General Relativity gives and a flow can carry;
\item tension as a pressure deficit, derived (Part 7);
\item the size of a spinning body's imprint (Part 16);
\item why each particle's knot has the mass it has, and the strength of charge relative to gravity (Part 17).
\end{itemize}

\section{How GPT Can Be Proved Wrong}

\argm{GPT is a testable account, and its tests are stated so that a result can count against it} (Part 29). One needs no new instrument, and the rest need only instruments that exist:
\begin{itemize}
\item \textbf{Pairs of stars, by direction.} GPT expects the relaxation boundary closer ahead of a moving mass than behind, so pairs of stars aligned with their motion should depart from Newton's law at a smaller separation. The Gaia data exist. The current analyses of wide binaries disagree with each other.\fn{chae2024}\textsuperscript{,}\fn{banik2024}
\item \textbf{A held body's weight and its structure.} The same pieces weighed in different arrangements, in vacuum, at matched height, on a mass comparator. The standard account expects a difference far below one part in $10^{15}$; a difference at one part in $10^{10}$ would lie beyond it.
\item \textbf{A crystal turned} on a torsion balance, watching its centre of gravity; \textbf{a spinning rotor}, watching for an imprint beyond its shape; \textbf{a clock} with its internal order changed.
\end{itemize}
\argm{Any of these can count against GPT, and so can a dark-matter particle found in the amounts the galaxies need,} or pairs of stars that follow Newton's law in every direction once GPT's account of the galaxy's flow is in hand.

\clearpage
\section*{The Restoration of Reality}
\addcontentsline{toc}{section}{The Restoration of Reality}

Every number stands. What was missing was the cause.

General Relativity says that nothing acts on a falling body, while buildings stand on their foundations, oceans rise beneath the Moon, galaxies turn, and every child who has ever fallen has been caught by the Earth. Graviton Pressure Theory says that something does act: a medium that fills space, flows into every mass, and presses on everything it meets.

We have named that medium, and set its inflow by Newton's law. We have placed the energy of gravity in its pressure, where General Relativity could not place it. We have followed the inflow from the apple to the planets, from the clock on the ground to the clock in orbit, from the bending of starlight to the edge of a black hole, and found the measured numbers at every step. We have built matter, time and order from the same medium, and we have stated, plainly, what we do not yet know.

This framework does not ask to be believed. It asks to be tested. Its tests are named, and most are within reach of instruments that already exist.

It is not the end of inquiry. It is a return to the question physics set aside: not only how much, but what. Causality is the standard; every claim must name what touches what. Truth is the measurement; every claim must pass it.

\vspace{1.5em}
\begin{center}\large
\textbf{This is Graviton Pressure Theory.}\\[0.6em]
\textbf{This is the restoration of reality.}\\[0.6em]
\textbf{And we are only just beginning.}
\end{center}

\clearpage
\section*{Changes in the Second Edition}
\addcontentsline{toc}{section}{Changes in the Second Edition}

This list records what changed from the first edition and why. Changes to the argument alter what the paper claims or how it supports a claim; clarifications improve wording, sources or presentation and leave the argument as it was.

\subsection*{Changes to the argument}
\begin{itemize}
\item \textbf{General Relativity is stated fairly.} The first edition said that General Relativity fails the test of physical realism, and that GPT replaces its foundation. This edition says that the two can agree on every number, that they cannot agree on what exists, and that GPT names the mechanism that the geometry leaves unnamed (Part 8).
\item \textbf{The claims are restated in the terms of the revised parts.} Fields as agents of structure become the medium's inflow, swirl and waves (Parts 16, 22 and 28). Gravity as directional field pressure becomes the inflow and its pressure. Mass as resistance to compression becomes energy slowed down, circulation held in a knot (Part 17). Time, force, inertia and motion as emergent become time as a structure's renewal (Part 18), every force with its carrier (Part 27), and momentum as the path a body carves through the medium (Part 11). Matter pressing rather than pulling becomes the rule that every force names its carrier (Part 7).
\item \textbf{Reconstructions withdrawn or restated:} the force equation as surface-dependent is withdrawn, since the delivery tracks a body's energy, not its surface; inertia as memory becomes the carved path; ``re-directional coherence'' and ``structural permission'' are replaced by the medium's inflow and pressure.
\item \textbf{Added:} the mechanism in one place, with the numbers it reproduces (Figure~\ref{fig:one}); the framework part by part (Table~\ref{tab:parts}); what was withdrawn; what is still owed; and how GPT can be proved wrong.
\end{itemize}

\subsection*{Clarifications}
\begin{itemize}
\item ``The final rupture'' and ``there is no going back to spacetime'' are restated: on what exists, only one account can be right, and the choice rests on mechanism. The closing lines are kept.
\item ``Testable, empirical, falsifiable'' is kept, with the tests named (Part 29).
\item Sources are added.
\end{itemize}

\bibliography{part30}
\end{document}
